\section{Related Work}
\label{sec:related}

\paragraph{Tail latency mitigation.}
Dean and Barroso~\cite{dean2013tail} established hedged requests and
tied requests as primary tools for tail latency in warehouse-scale
systems. Subsequent work refined these ideas: Vulimiri et
al.~\cite{vulimiri2013low} analyzed redundancy strategies under various
latency distributions; Ananthanarayanan et
al.~\cite{ananthanarayanan2013effective} applied speculation to
stragglers in data analytics. Our work shows that these reactive
mechanisms, designed for random stragglers, are ineffective against
\emph{persistent} slow faults---a distinct failure class requiring
proactive detection.

\paragraph{Gray failure detection.}
Huang et al.~\cite{huang2017gray} characterized gray failures as
differential observability problems and proposed detection via
cross-component correlation. Gunawi et al.~\cite{gunawi2018fail}
cataloged failure modes in cloud systems, finding that performance
degradation accounts for a significant fraction. Our departure-interval
detection builds on this work by providing a load-invariant signal that
distinguishes genuine slowdown from queueing-induced latency increases.

\paragraph{Outlier detection and blacklisting.}
Cassandra's dynamic snitch~\cite{lakshman2010cassandra} uses
exponentially decaying latency scores to route away from slow replicas.
Netflix's adaptive routing~\cite{netflix2020zuul} extends this with
real-time outlier detection. We show that blacklisting is effective at
moderate load but exhibits catastrophic oscillation at high load due to
the capacity cost of node removal.

\paragraph{Load balancing theory.}
Mitzenmacher~\cite{mitzenmacher2001power} proved the exponential
improvement of P2C over random routing for homogeneous servers.
Join-Shortest-Queue (JSQ), also known as Least-Outstanding-Requests,
is optimal for homogeneous systems~\cite{winston1977optimality} but
requires global queue-state information. Our evaluation
(Section~\ref{sec:eval-lor}) shows that JSQ's deterministic routing
creates a synchronization failure mode under fault isolation: when
faulty nodes are removed, all arrivals concentrate on the same
shortest-queue worker, causing herd behavior that P2C's randomized
sampling avoids. We extend the P2C analysis to \emph{heterogeneous}
servers (faulty nodes with different service rates), deriving the
capacity-proportional traffic distribution that underlies the severity
non-monotonicity result.

\paragraph{Capacity planning under failures.}
The $N+k$ redundancy model is standard practice, but the relationship
between load factor, fault fraction, and per-node utilization ceiling
has not been formalized for slow faults specifically. Our operating
envelope formula $\fmax = 1 - L/U$ provides a simple, validated tool
for this purpose.

\paragraph{Adaptive load shedding.}
DAGOR~\cite{zhang2021dagor} and CoDel~\cite{nichols2012codel} address
overload through adaptive shedding, but focus on aggregate system load
rather than per-node fault isolation. Our results show that under
queue-aware routing, the specific response strategy (graduated
shedding vs.\ full isolation) is inconsequential ($< 2$\,ms P99
difference); detection speed is the key bottleneck. This contrasts
with the above works' emphasis on response strategy design.
